\documentclass[10pt,a4paper]{article}
\usepackage[utf8]{inputenc}
\usepackage{amsmath}
\usepackage{amsfonts}
\usepackage{amssymb}
\usepackage{amsthm}
\usepackage{hyperref}

\usepackage{multicol}
\usepackage{fancyhdr}
\usepackage[inline]{enumitem}
\usepackage{tikz}
\usepackage{tikz-cd}
\usetikzlibrary{calc}
\usetikzlibrary{shapes.geometric}
\usepackage[margin=0.5in]{geometry}
\usepackage{xcolor}
\usepackage{ esint }

\hypersetup{
    colorlinks=true,
    linkcolor=blue,
    filecolor=magenta,      
    urlcolor=cyan,
    pdftitle={Tensors},
    pdfpagemode=FullScreen,
    }

%\urlstyle{same}

\newcommand{\CLASSNAME}{Math 5315 -- Numerical Methods: Partial Differential Equations}
\newcommand{\STUDENTNAME}{Paul Carmody}
\newcommand{\ASSIGNMENT}{NOTES }
\newcommand{\DUEDATE}{December 17, 2025}
\newcommand{\SEMESTER}{Fall 2025}
\newcommand{\SCHEDULE}{TTh 2:00 -- 3:15}
\newcommand{\ROOM}{Crawford 110}

\newcommand{\MMN}{M_{m\times n}}
\newcommand{\FF}{\mathcal{F}}

\pagestyle{fancy}
\fancyhf{}
\chead{ \fancyplain{}{\CLASSNAME} }
%\chead{ \fancyplain{}{\STUDENTNAME} }
\rhead{\thepage}
\input{../MyMacros}

\newcommand{\RED}[1]{\textcolor{red}{#1}}
\newcommand{\BLUE}[1]{\textcolor{blue}{#1}}

\begin{document}
\begin{center}
	\textbf{\Large{High Order Weighted
Essentially Nonoscillatory
Schemes for \\Convection
Dominated Problems}}
\end{center}
\begin{center}
	A review by Paul Carmody\\Math 5315 Numerical Methods for PDE: Fall 2025
\end{center}

\HLINE
\begin{abstract}
	Provide a synapses and understanding of the article ``High Order Weighted
Essentially Nonoscillatory
Schemes for Convection
Dominated Problems" by Chi-Wang Shu.  The background and motivation for the WENO process is described along with its foundation in ENO (and its limitations).  The WENO methodology is then described in a step by step manner.  We measure than the behavior of the non-linear weights (decribed below).  This article ends with a discussion measuring the merits, advantages and disadvantages of this approach.
\end{abstract}

\HLINE
The Weighted Essentially Non-Oscillatory (WENO) scheme is a high-resolution numerical method widely used for solving hyperbolic partial differential equations (PDEs), particularly those involving discontinuities, shocks, or sharp gradients.\cite{weno} It is designed to achieve high-order accuracy in smooth regions while simultaneously suppressing spurious oscillations near discontinuities, a common problem with standard high-order schemes (like central difference schemes).

Here is a detailed explanation of the WENO scheme, focusing on its application in numerical methods for solving PDEs.

\begin{description}

\item Reasoning \& Explanation

\begin{description}

\item  \textbf{Background and Motivation}

Hyperbolic conservation laws often take the form:
$$\frac{\partial u}{\partial t} + \frac{\partial f(u)}{\partial x} = 0$$
where $u$ is the conserved quantity and $f(u)$ is the flux function.

When solving these equations numerically, especially using finite difference or finite volume methods, achieving high accuracy (order $p \ge 3$) is desirable. However, Godunov's theorem (or the TVD theorem) states that linear schemes that are Total Variation Diminishing (TVD) must be at most first-order accurate near extrema. High-order linear schemes (like central differences) introduce oscillations (Gibbs phenomena) near discontinuities.

Non-linear schemes, such as TVD schemes (e.g., MUSCL), overcome this limitation but often cap the accuracy at second order. The WENO scheme is a non-linear approach that achieves high-order accuracy (typically fifth order) in smooth regions while maintaining non-oscillatory properties near shocks.

\item \textbf{The ENO Foundation}

WENO is an improvement upon the Essentially Non-Oscillatory (ENO) scheme, developed by Harten, Engquist, Osher, and Chakravarthy.

ENO Principle: The ENO scheme achieves non-oscillatory behavior by adaptively choosing the smoothest stencil (a set of neighboring grid points) from a collection of candidate stencils to construct the numerical flux or interpolation polynomial. If a stencil spans a discontinuity, it is discarded in favor of a smoother one.\cite{}

Limitation of ENO: ENO is only "essentially" non-oscillatory, meaning the error is $O(\Delta x^p)$ in smooth regions, but the scheme is not optimally accurate. The selection process is discrete (either choose stencil A or stencil B), which leads to a non-smooth flux function and a slight reduction in convergence rate compared to its potential order.

\item \textbf{The WENO Methodology}

The WENO scheme, introduced by Liu, Osher, and Chan (WENO-L), and later refined by Jiang and Shu (WENO-JS), overcomes the limitations of ENO by using a \textit{convex combination} of all candidate stencils, weighted by non-linear weights.\cite{eno}

\begin{enumerate}

\item \textbf{Flux Splitting (Required for Hyperbolic Systems)}

To ensure stability and handle the direction of information propagation, the flux function $f(u)$ is typically split into positive and negative parts based on the sign of the eigenvalues of the Jacobian $\partial f / \partial u$. This is often done using the Lax-Friedrichs splitting or characteristic decomposition:
$$f(u) = f^+(u) + f^-(u)$$
The numerical flux $\hat{f}_{i+1/2}$ is then calculated as the sum of the fluxes derived from $f^+$ and $f^-$. We focus on calculating the numerical flux for $f^+$ (the procedure is symmetric for $f^-$).

\item \textbf{Candidate Stencils and Interpolation}

For a $p$-th order WENO scheme (e.g., $p=5$), we define $r$ small stencils, $S_k$, each containing $r$ points, such that their union forms a large stencil $S$ of size $2r-1$. For 5th-order WENO ($r=3$), we use three 3-point stencils:

| Stencil Index ($k$) | Stencil Points |
|---|---|
| 0 | $u_{i-2}, u_{i-1}, u_i$ |
| 1 | $u_{i-1}, u_i, u_{i+1}$ |
| 2 | $u_i, u_{i+1}, u_{i+2}$ |

For each small stencil $S_k$, a high-order polynomial interpolation $P_k(x)$ is constructed to approximate the flux function $f(u)$ at the interface $x_{i+1/2}$. This yields a candidate numerical flux $\hat{f}_{i+1/2}^{(k)}$.

\item  \textbf{Linear Weights (Ideal Weights)}

If the solution $u$ were perfectly smooth, we could combine these candidate fluxes using fixed, pre-determined linear weights, $d_k$, to achieve the desired high order (e.g., 5th order). These weights satisfy $\sum_{k} d_k = 1$.

For 5th-order WENO-JS, the linear weights are:
$$d_0 = 0.1, \quad d_1 = 0.6, \quad d_2 = 0.3$$

The high-order flux approximation $\hat{f}_{i+1/2}^{HO}$ would be:
$$\hat{f}_{i+1/2}^{HO} = \sum_{k=0}^{r-1} d_k \hat{f}_{i+1/2}^{(k)}$$

\item \textbf{ Smoothness Indicators ($\beta_k$)}\cite{beta}

The core of the WENO scheme is the determination of how smooth the solution is within each candidate stencil $S_k$. This is quantified by a smoothness indicator, $\beta_k$. $\beta_k$ is typically a measure of the $L_2$ norm of the derivatives of the interpolation polynomial $P_k(x)$ over the stencil.

A common definition for $\beta_k$ (for 5th-order WENO-JS) involves quadratic forms of the cell averages:
$$\beta_k = \sum_{l=1}^{r-1} \Delta x^{2l-1} \int_{x_{i-1/2}}^{x_{i+1/2}} \left(\frac{d^l P_k(x)}{dx^l}\right)^2 dx$$
In practice, $\beta_k$ is calculated using discrete differences of the cell values $u_i$. A smaller $\beta_k$ indicates a smoother solution within stencil $S_k$. If a discontinuity lies within $S_k$, $\beta_k$ will be large.

\item \textbf{Non-linear Weights ($\omega_k$)}

The non-linear weights $\omega_k$ are calculated based on the linear weights $d_k$ and the smoothness indicators $\beta_k$. The goal is to assign a weight close to $d_k$ if the stencil is smooth, and a weight close to zero if the stencil contains a discontinuity.

The standard WENO-JS formulation uses:
$$\alpha_k = \frac{d_k}{(\epsilon + \beta_k)^2}$$
$$\omega_k = \frac{\alpha_k}{\sum_{j=0}^{r-1} \alpha_j}$$
where $\epsilon$ is a small positive number (e.g., $10^{-6}$) introduced to prevent division by zero when $\beta_k$ is zero and to stabilize the scheme.

Behavior of Non-linear Weights:
1.  Smooth Region: If all stencils are smooth, $\beta_k \approx O(\Delta x^2)$. The weights $\omega_k$ approach the linear weights $d_k$, and the scheme achieves the full high order (e.g., 5th order).
2.  Discontinuity Present: If stencil $S_j$ contains a discontinuity, $\beta_j$ will be large, causing $\alpha_j$ and consequently $\omega_j$ to approach zero. The scheme effectively switches to a lower-order, non-oscillatory interpolation based on the remaining smooth stencils.

\item  \textbf{Final Numerical Flux}

The final numerical flux $\hat{f}_{i+1/2}$ is the weighted average of the candidate fluxes:
$$\hat{f}_{i+1/2} = \sum_{k=0}^{r-1} \omega_k \hat{f}_{i+1/2}^{(k)}$$

This flux is then used in the semi-discrete finite volume formulation:
$$\frac{d u_i}{d t} = -\frac{1}{\Delta x} (\hat{f}_{i+1/2} - \hat{f}_{i-1/2})$$

\end{enumerate}
\item \textbf{ Advantages and Disadvantages}

\begin{description}

\item \textbf{High Accuracy:} Achieves high order (typically 5th or 7th) in smooth regions.

\item \textbf{Non-Oscillatory} Effectively suppresses spurious oscillations near shocks and discontinuities. 

\item \textbf{Robustness} Highly robust for complex flow problems, including turbulence and compressible flows. 

\item \textbf{Complexity} Significantly more complex to implement than linear schemes due to flux splitting, multiple stencil calculations, and non-linear weighting. |

\item \textbf{Computational Cost} Higher computational cost per time step compared to lower-order schemes (e.g., 2nd order TVD) due to the multiple flux evaluations and weight calculations. 

\item \textbf{Accuracy Near Extrema} While high-order in smooth regions, the scheme typically drops to 3rd order accuracy near critical points (smooth extrema) due to the non-linear weighting function. (This led to the development of improved schemes like WENO-M and WENO-Z). 

\item \textbf{Summary} 

WENO achieves high order in smooth regions, suppresses oscillations near shocks, and remains robust for complex flows, while being more computationally demanding. Improvements such as WENO-M and WENO-Z enhance accuracy near extrema.

\end{description}

\item \textbf{Variants of WENO}\cite{current}

\begin{enumerate}
\item \textbf{WENO-JS (Jiang and Shu):} 

The standard, widely used formulation described above.
\item \textbf{WENO-M (Modified):} 

A modification that aims to restore the optimal order of accuracy (5th order) even at smooth extrema by modifying the non-linear weights.
\item \textbf{WENO-Z (Zhong):} 

Uses a different definition for the non-linear weights that incorporates a global smoothness indicator, often leading to better performance and higher accuracy near critical points than WENO-JS.

\end{enumerate}

\end{description}
---

\item Final Answer Summary

The Weighted Essentially Non-Oscillatory (WENO) scheme is a high-resolution numerical method for solving hyperbolic PDEs. It achieves high-order accuracy (e.g., 5th order) in smooth regions while maintaining non-oscillatory behavior near discontinuities (shocks).

The core mechanism involves:
\begin{enumerate}

\item  \textbf{Flux Splitting: }

Decomposing the flux into positive and negative components.
\item \textbf{Multiple Stencils: }

Constructing candidate numerical fluxes from several small, overlapping stencils.
\item  \textbf{Smoothness Indicators ($\beta_k$): }

Quantifying the smoothness of the solution within each stencil.
\item \textbf{Non-linear Weights ($\omega_k$): }

Combining the candidate fluxes using weights that adaptively favor the smoothest stencils. If a stencil contains a discontinuity (high $\beta_k$), its weight $\omega_k$ approaches zero.

\end{enumerate}
The final numerical flux is a convex combination of the candidate fluxes:
$$\hat{f}_{i+1/2} = \sum_{k} \omega_k \hat{f}_{i+1/2}^{(k)}$$

\end{description}
\textbf{Conclusion:}

The Weighted Essentially Non-Oscillatory (WENO) scheme is a high-resolution numerical method for solving hyperbolic PDEs that uses adaptive, non-linear weights with multiple stencils to achieve high-order accuracy in smooth regions and suppress oscillations near discontinuities. Its core process involves flux splitting, candidate stencil selection, smoothness indicators, and convex combination of candidate fluxes using non-linear weights, resulting in a final numerical flux of \(\hat{f}_{i+1/2} = \sum_{k} \omega_k \hat{f}_{i+1/2}^{(k)}\).


Answer: The Weighted Essentially Non-Oscillatory (WENO) scheme is a high-resolution numerical method for solving hyperbolic PDEs that uses adaptive, non-linear weights with multiple stencils to achieve high-order accuracy in smooth regions and suppress oscillations near discontinuities. Its core process involves flux splitting, candidate stencil selection, smoothness indicators, and convex combination of candidate fluxes using non-linear weights, resulting in a final numerical flux of \(\hat{f}_{i+1/2} = \sum_{k} \omega_k \hat{f}_{i+1/2}^{(k)}\).

\begin{thebibliography}{4}
\bibitem{weno}
Liu, X.-D., Osher, S., and Chan, T. (1994). "Weighted Essentially Non-Oscillatory Schemes." *Journal of Computational Physics*, 115(1), 200-212.

\bibitem{impl}
Jiang, G.-S., and Shu, C.-W. (1996). "Efficient Implementation of Weighted ENO Schemes." *Journal of Computational Physics*, 126(1), 202-228.

\bibitem{beta}
 
 Shu, C.-W. (1998). "Essentially Non-Oscillatory and Weighted Essentially Non-Oscillatory Schemes for Hyperbolic Conservation Laws." In *Advanced Numerical Approximation of Nonlinear Hyperbolic Equations* (pp. 325-432). Springer, Berlin, Heidelberg.

\bibitem{current} 

Shu, C.-W. (2009). "High Order Weighted Essentially Non-Oscillatory Schemes for Convection Dominated Problems." *SIAM Review*, 51(1), 82-126.

\end{thebibliography}

\end{document}